\documentclass[11pt]{article}
\usepackage[margin=2.54cm]{geometry}
\usepackage[dvipsnames,usenames]{color}
\usepackage{amsfonts,amsmath,amssymb,amsthm,mathtools}
\usepackage{paralist}
\usepackage{algorithmic,algorithm}
\usepackage{bm}
\usepackage{xspace}
\usepackage[pagebackref,letterpaper=true,colorlinks=true,pdfpagemode=none,urlcolor=blue,linkcolor=blue,citecolor=BrickRed,pdfstartview=FitH]{hyperref}
\usepackage{xspace,prettyref}
\usepackage{color}
%\usepackage{MnSymbol}
%% For hyperlinks. Should always be the last package.
%\usepackage[colorlinks,urlcolor=blue,citecolor=blue,linkcolor=blue]{hyperref}
\let\pref=\prettyref
\newcommand{\savehyperref}[2]{\texorpdfstring{\hyperref[#1]{#2}}{#2}}
\newcommand{\comment}[1]{ {\color{BrickRed} \footnotesize[#1]}\marginpar{\footnotesize\textbf{\color{red} To Do!}}}
\newcommand{\ignore}[1]{}

\newtheorem{theorem}{Theorem}
\newtheorem{lemmata}[theorem]{Lemmata}
\newtheorem{lemma}[theorem]{Lemma}
\newtheorem{claim}[theorem]{Claim}
\newtheorem{subclaim}{Subclaim}
\newtheorem{proposition}[theorem]{Proposition}
\newtheorem{corollary}[theorem]{Corollary}
\newtheorem{fact}[theorem]{Fact}
\newtheorem{conjecture}[theorem]{Conjecture}
\newtheorem{question}[theorem]{Question}
\newtheorem{example}[theorem]{Example}
\newtheorem{definition}{Definition}
\theoremstyle{definition}
\newtheorem{remark}{Remark}
\newtheorem{observation}{Observation}

\def\R{\RR}
\def\RR{\mathbb R}
\def\ZZ{\mathbb Z}
\def\eqdef{~\triangleq~}
\def\eps{\varepsilon}
\def\bar{\overline}
\def\floor#1{\lfloor {#1} \rfloor}
\def\ceil#1{\lceil {#1} \rceil}
\def\script#1{\mathcal{#1}}

\newenvironment{proofof}[1]{\smallskip\noindent{\bf Proof of #1:}}%
        {\hspace*{\fill}$\Box$\par}

%% Hypercube related macros

\def\B{{\mathsf D}}
\def\Hyp{{\sf Hyp}}
\def\D{{\sf K}}
\def\G{{\mathsf G}}
\def\d{{\sf d}}
\def\a{{\mathfrak a}}
\def\b{{\mathfrak b}}
\def\c{{\mathfrak c}}
\def\db{{\mathfrak d}}
\def\e{{\mathfrak e}}
\def\f{{\mathfrak f}}
\def\h{{\mathfrak h}}
\def\g{{\mathfrak g}}

\newcommand{\PTIME}{\mathbb{P}}


\newcommand{\hu}{\widehat{u}}

\newcommand{\hcd}{\mathsf{hcd}}
\newcommand{\VG}{VG}
\newcommand{\md}[1]{\ (\operatorname{mod} #1)}

\newcommand{\cupdot}{\sqcup}

%% Calligraphic letters

\newcommand{\cA}{{\cal A}}
\newcommand{\cB}{{\cal B}}
\newcommand{\cC}{{\cal C}}
\newcommand{\cD}{{\cal D}}
\newcommand{\cE}{{\cal E}}
\newcommand{\cF}{{\cal F}}
\newcommand{\cG}{{\cal G}}
\newcommand{\cH}{{\cal H}}
\newcommand{\cI}{{\cal I}}
\newcommand{\cJ}{{\cal J}}
\newcommand{\cL}{{\cal L}}
\newcommand{\cM}{{\cal M}}
\newcommand{\cP}{{\cal P}}
\newcommand{\cR}{{\cal R}}
\newcommand{\cS}{{\cal S}}
\newcommand{\cT}{{\cal T}}
\newcommand{\cU}{{\cal U}}
\newcommand{\cV}{{\cal V}}

%% Bold letters

\newcommand{\bS}{{\bf S}}

\newcommand{\NN}{\mathbb{N}}

%% Hyper-linked References
\newcommand{\Sec}[1]{\hyperref[sec:#1]{\S\ref*{sec:#1}}} %section
\newcommand{\Eqn}[1]{\hyperref[eq:#1]{(\ref*{eq:#1})}} %equation
\newcommand{\Fig}[1]{\hyperref[fig:#1]{Fig.\,\ref*{fig:#1}}} %figure
\newcommand{\Tab}[1]{\hyperref[tab:#1]{Tab.\,\ref*{tab:#1}}} %table
\newcommand{\Thm}[1]{\hyperref[thm:#1]{Theorem\,\ref*{thm:#1}}} %theorem
\newcommand{\Lem}[1]{\hyperref[lem:#1]{Lemma\,\ref*{lem:#1}}} %lemma
\newcommand{\Prop}[1]{\hyperref[prop:#1]{Prop.~\ref*{prop:#1}}} %property
\newcommand{\Cor}[1]{\hyperref[cor:#1]{Corollary~\ref*{cor:#1}}} %corollary
\newcommand{\Def}[1]{\hyperref[def:#1]{Definition~\ref*{def:#1}}} %definition
\newcommand{\Alg}[1]{\hyperref[alg:#1]{Alg.~\ref*{alg:#1}}} %algorithm
\newcommand{\Ex}[1]{\hyperref[ex:#1]{Ex.~\ref*{ex:#1}}} %example
\newcommand{\Clm}[1]{\hyperref[clm:#1]{Claim~\ref*{clm:#1}}} %example


\def\yes{{\sf YES }}
\def\no{{\sf NO }}
\def\Lip{{\sf Lip}}
\def\poly{{\sf poly}}

\newcommand{\ab}{\hbox{ab}}
\newcommand{\abs}[1]{{|#1|}}
\newcommand{\hd}[2]{{||#1 - #2||_1}}
\newcommand{\proj}{{\mathsf{proj}}}
\newcommand{\drop}{\hbox{drop}}
\newcommand{\dist}[2]{\eps_{#1,#2}}
\newcommand{\hM}{\widehat{M}}
\newcommand{\hX}{\widehat{X}}

\newcommand{\flo}[1]{\lfloor #1 \rfloor}

\newcommand{\EXPTIME}{\mathbb{EXP}}
\newcommand{\NP}{\mathbb{NP}}
\newcommand{\coNP}{\textrm{co-}\mathbb{NP}}
\newcommand{\NEXP}{\mathbb{NEXP}}



\newcommand{\Sesh}[1]{{\color{red} Sesh says: #1}}

\title{CSE204: Exercises in NP-completeness\\
{\small \copyright C. Seshadhri, 2025}}

\date{}

\begin{document}

\maketitle

Please cite any collaborations
or online sources. All assignments must be done in Latex and submitted as pdfs.

\section{Easier problems}

Try to make some attempts by yourself, before collaborating with others. Ideally, do not use GenAI for these
problems. The answers are not that long. We can use any previous statement to prove subsequent ones.

\begin{enumerate}

    \item Prove that the language

    $\cL = \{\langle M, x, 1^t\rangle | \textrm{M is a non-determinisitic TM
     that runs in t steps and accepts $x$}\}$ is $\NP$-complete.

    \item Consider the Cook-Levin reduction for a language in $\cL \in \PTIME$. It will still construct a Boolean formula
        $\phi_x$, given an input $x$, such that $x \in \cL$ iff $\phi_x$ is satisfiable. Prove that a satisfying assignment
        for $\phi_x$ can be computed in polynomial time.

    \item Complete the proof that CLIQUE is $\NP$-complete. (Show that the reduction given in class works.)

    \item Prove that 4SAT is $\NP$-complete. 

    \item Prove formally that Integer Linear Programming is $\NP$-complete.

    \item $\cL$ is $\coNP$-complete if: $\cL \in \coNP$ and $\forall \cL' \in \coNP$, $\cL' \leq_p \cL$.
    Prove that the class of $\coNP$-complete problems is exactly the set of co-($\NP$-complete) problems.

\end{enumerate}


\section{ChatGPT questions} 

Put this question ChatGPT and get the answer. You are allowed to discuss further. Please put a link to the response and your conversation. 
Given the response, you must first say whether: (i) the response is correct, (ii) the response is false, or (iii) you are really not sure.
In the first case, give a succinct explanation of the proof. In the second case, explain why it is false.
In the third case, explain your confusion. You do not need to write too much.

\begin{enumerate}
    \item NAE-3SAT is the language of 3-CNFs with an assignment that does not make all literals in a clause have the same value.
    Prove that NAE-3SAT is $\NP$-complete. 
    \item Prove that if FACTORING is $\coNP$-complete, then $TAUT \leq_{p(n)} SAT$ (recall
    that $TAUT$ is the language of tautologically true Boolean formulas.)
\end{enumerate}


\section{More challenging exercises}

\begin{enumerate}
    \item Consider a simple, undirected graph $G$. A \emph{cut} is a partition of the vertices in two sets $S, T$. The value of the cut
    is the number of edges ``crossing" the cut, i.e. the number of edges between $S$ and $T$. Define the language 
        MAXCUT $= \{\langle G, k\rangle | \textrm{$G$ has a cut of value at least $k$} \}$. Prove that MAXCUT is $\NP$-complete.

    (Hint: A more tricky reduction. Prove NAESAT $\leq_p$ MAXCUT. Start with $\phi$. For each variable, create a sufficiently large bi-clique. So create a bunch
    of nodes labeled $x_i$ and a bunch labeled $\overline{x_i}$, and connect all pairs of nodes with $x_i$ to those with $\overline{x_i}$. For each clause $y_i \vee y_j \vee y_k$,
    create a triangle between nodes with these labels. Work out the details to get the reduction correct.)
    \item Prove Berman's theorem, stated as follows. A language is called \emph{unary} if every string in the language is of the
    form $1^i$, for some $i \in \NN$. If a unary language is $\NP$-complete, then $\PTIME = \NP$.

    (Hint: start with Boolean formula $\phi$. Using the reduction to a unary language, we get a unary string $s_\phi$.
    Take a variable in $\phi$, set it to both $0$ and $1$, to get two formulas $\phi_0$ and $\phi_1$. Observe that $\phi$
    is satisfiable iff either $\phi_0$ or $\phi_1$ is satisfiable. Apply the unary reduction to these new formulas.
    Keep repeating this process with $\phi_0$ and $\phi_1$, accumulating more unary strings. Show that this
    process can be run to termination in polynomial time. Use the unary strings to ``prune" branches of this process.)



\end{enumerate}
\end{document}
